% Exported from Math++. Compile with pdflatex.
\documentclass[11pt,a4paper]{article}
\usepackage[T1]{fontenc}
\usepackage[utf8]{inputenc}
\usepackage{lmodern}
\usepackage{amsmath,amssymb,amsthm,mathtools}
\usepackage{graphicx,booktabs,xcolor}
\usepackage[hidelinks]{hyperref}
\usepackage[margin=25mm]{geometry}
\theoremstyle{plain}
\newtheorem{theorem}{Theorem}
\newtheorem{lemma}[theorem]{Lemma}
\newtheorem{corollary}[theorem]{Corollary}
\newtheorem{proposition}[theorem]{Proposition}
\theoremstyle{definition}
\newtheorem{definition}[theorem]{Definition}
\newtheorem{example}[theorem]{Example}
\theoremstyle{remark}
\newtheorem*{remark}{Remark}
\setlength{\parindent}{0pt}\setlength{\parskip}{0.6em}
\title{Auto-evaluation}
\date{}
\begin{document}
\maketitle
Type the maths and the answer appears after it in grey, without \texttt{= ?}. The answers are only
shown: the file keeps what you typed. \texttt{Ctrl+Enter} on a line writes its answer into the file.

\section*{Integrals}

\begin{equation*}
  \int x^2 \, dx
\end{equation*}
\begin{equation*}
  \int_0^1 x^2 \, dx
\end{equation*}
\begin{equation*}
  \int_0^1 \int_0^2 x y \, dy \, dx
\end{equation*}
\begin{equation*}
  \iint_{[0, 1] \times [0, 2]} x y \, dA
\end{equation*}
\begin{equation*}
  \iiint_{[0, 1] \times [0, 1] \times [0, 1]} x y z \, dV
\end{equation*}
A region written as a condition switches coordinates on its own: polar for a disc,
spherical for a ball.

\begin{equation*}
  \iint_{x^2 + y^2 \le 4} (x^2 + y^2) \, dA
\end{equation*}
\begin{equation*}
  \iiint_{x^2 + y^2 + z^2 \le 1} 1 \, dV
\end{equation*}
Line and surface integrals around a circle or over a sphere:

\begin{equation*}
  \oint_{x^2 + y^2 = 1} (-y, x) \cdot d\mathbf{r}
\end{equation*}
\begin{equation*}
  \oiint_{x^2 + y^2 + z^2 = 1} (x, y, z) \cdot d\mathbf{S}
\end{equation*}
\section*{Derivatives}

\begin{equation*}
  f(x) \coloneqq  x^3 \sin x
\end{equation*}
\begin{equation*}
  \frac{d}{dx} x^3
\end{equation*}
\begin{equation*}
  f'(x)
\end{equation*}
\begin{equation*}
  f''(0)
\end{equation*}
\begin{equation*}
  \partial_x (x^2 y)
\end{equation*}
\begin{equation*}
  D_x \ln x
\end{equation*}
\begin{equation*}
  \frac{\partial^2}{\partial x \partial y} x^2 y^3
\end{equation*}
\begin{equation*}
  \nabla (x^2 y)
\end{equation*}
\begin{equation*}
  \nabla \cdot (x^2, y^2, z^2)
\end{equation*}
\begin{equation*}
  \nabla \times (-y, x, 0)
\end{equation*}
\begin{equation*}
  \nabla^2 (x^2 + y^2 + z^2)
\end{equation*}
\begin{equation*}
  \frac{\partial(u + v, u - v)}{\partial(u, v)}
\end{equation*}
\begin{equation*}
  g(x, y) \coloneqq  x^2 y
\end{equation*}
\begin{equation*}
  dg
\end{equation*}
\section*{Limits, sums, products and series}

\begin{equation*}
  \lim_{x \to 0} \frac{\sin x}{x}
\end{equation*}
\begin{equation*}
  \sum_{n=1}^{\infty} \frac{1}{n^2}
\end{equation*}
\begin{equation*}
  \prod_{k=1}^{5} k
\end{equation*}
\begin{equation*}
  \operatorname{series}_{x=0}^{4} e^x
\end{equation*}
\section*{Transforms}

\begin{equation*}
  \mathcal{L}\{t^2 e^{-t}\}
\end{equation*}
\begin{equation*}
  \mathcal{L}_{t \to s}\left\{ \sin 2t \right\}
\end{equation*}
\begin{equation*}
  \mathcal{L}^{-1}\left\{ \frac{1}{s^2 + 4} \right\}
\end{equation*}
\begin{equation*}
  \mathcal{F}\{e^{-x^2}\}
\end{equation*}
\begin{equation*}
  \mathcal{F}^{-1}\left\{ \frac{2}{1 + \omega^2} \right\}
\end{equation*}
\section*{Equations}

An equation with an unknown is solved; a system can continue after a comma.

\begin{equation*}
  x^2 - 5x + 6 = 0
\end{equation*}
\begin{equation*}
  x^2 - 4 < 0
\end{equation*}
\begin{equation*}
  x + y = 3, \quad x - y = 1
\end{equation*}
\begin{equation*}
  \nabla (x y) = \lambda \nabla (x^2 + y^2 - 1), \quad x^2 + y^2 = 1
\end{equation*}
A differential equation shows its solution:

\begin{equation*}
  y'' + y = 0, \quad y(0) = 1, y'(0) = 0
\end{equation*}
\begin{equation*}
  \frac{dy}{dx} = x y
\end{equation*}
\section*{Matrices}

\begin{align*}
  \det \begin{pmatrix} 1 & 2 \\ 3 & 4 \end{pmatrix}
\end{align*}
\begin{align*}
  \begin{pmatrix} 1 & 2 \\ 3 & 4 \end{pmatrix}^{-1}
\end{align*}
\begin{align*}
  \operatorname{eig} \begin{pmatrix} 2 & 1 \\ 1 & 2 \end{pmatrix}
\end{align*}
\section*{No answer shown}

Definitions, checks, expressions with nothing to work out, and blocks marked \texttt{\%noeval}:

\begin{equation*}
  a \coloneqq  3
\end{equation*}
\begin{equation*}
  \sin^2 x + \cos^2 x \overset{\checkmark }{=} 1
\end{equation*}
\begin{equation*}
  x^2 + 1
\end{equation*}
\begin{equation*}
  y = x^2 + 1
\end{equation*}
\begin{equation*}
  \int_0^1 x \, dx
\end{equation*}
\end{document}
